\section{Requirements-Based Business Onboarding and Progressive Capability Activation}
\label{sec:aion-requirements-onboarding}

\subsection{Decision: Do Not Classify the Customer by Business Size}

AION must not require a customer to identify as a ``small'', ``medium'' or ``large'' business during
onboarding.  These labels vary between countries, regulators and industries, and employee count or
turnover rarely describes the operational, security or intelligence capabilities that the customer
actually needs.  A one-person regulated consultancy may require stronger controls than an
unregulated company employing hundreds of people.

The onboarding process therefore constructs a requirements profile rather than assigning an
arbitrary company-size category.  This profile determines which capabilities AION recommends and
which parts of the interface it progressively reveals.  The customer's permanent brain, identity,
Business Map and data ownership remain the same regardless of the selected capabilities or
commercial package.

\subsection{Requirements Qualification}

The initial business conversation should establish only the information necessary to configure the
customer's operating environment.  The qualification process covers:

\begin{itemize}
  \item legal entities, trading names and brands;
  \item physical locations, operating countries, regions and currencies;
  \item divisions, departments, product families and business units;
  \item the number and types of users, including employees, contractors, advisers and guests;
  \item reporting structures, matrix relationships, projects and temporary teams;
  \item approval levels, delegated authority and separation-of-duty requirements;
  \item identity systems such as Microsoft Entra ID, Google Workspace, local identity or another
        approved single-sign-on provider;
  \item data classification, residency, retention and regulatory requirements;
  \item existing business systems, files, communications, calendars, CRM, accounting and databases;
  \item local, customer-server, private-cloud and optional external model requirements;
  \item audit, backup, disaster-recovery, availability and service-level expectations; and
  \item the customer's preferred balance between simplicity, control, cost and technical management.
\end{itemize}

Questions should be progressive.  AION should not present enterprise configuration merely because
it exists.  It should ask a deeper question only when a previous answer reveals a relevant need.  A
customer with one entity, one location and one owner should be able to complete onboarding without
encountering regional residency, matrix reporting or multi-entity consolidation settings.

\subsection{Operating-Complexity Profile}

Internally, AION may represent the customer's requirements as a composable operating-complexity
profile.  The profile is not a status label and must not be shown as a judgment about company size.
An illustrative profile is:

\begin{quote}
\texttt{Core business} \\
$+$ \texttt{Project operations} \\
$+$ \texttt{Multi-location} \\
$+$ \texttt{Multi-entity} \\
$+$ \texttt{Delegated authority} \\
$+$ \texttt{Enterprise identity} \\
$+$ \texttt{Regulated data} \\
$+$ \texttt{Private compute}
\end{quote}

Each component is independently activated, reviewed and removable.  The profile should be derived
from recorded requirements and approved by the customer.  AION must explain why a component was
recommended and must not treat the recommendation as commercial entitlement or consent.

\subsection{Progressive Disclosure of the Boardroom}

The same Boardroom adapts to the customer's requirements rather than splitting into unrelated
products:

\begin{itemize}
  \item A sole trader initially sees a simple Boardroom, Pilot, tasks, calendar, files and essential
        business controls.
  \item A project-based company receives project teams, budgets, milestones, risks, approvals and
        project reporting.
  \item A multi-site organisation receives location hierarchies, regional dashboards, local
        responsibility and geographically scoped permissions.
  \item A corporate group receives multiple legal entities, brands, currencies, inter-company
        boundaries and consolidated reporting.
  \item A regulated organisation receives explicit data residency, retention, legal hold, audit,
        stronger approval and private-compute controls.
  \item A complex organisation may add arbitrary divisions, matrix teams, temporary programmes and
        roles that do not map cleanly to traditional departments such as Sales or Finance.
\end{itemize}

Advanced structures remain available without cluttering the experience of a simpler business.  If
the customer's needs change, AION updates the requirements profile and proposes the additional
module; it does not require the customer to migrate to a different brain or rebuild the Business Map.

\subsection{Commercial Packaging Is a Separate Decision}

Organisational complexity and commercial packaging must remain independent.  A requirements
profile describes what the business needs.  An entitlement describes which supported capabilities
the customer has chosen to activate.  The initial commercial structure may be expressed as:

\begin{description}
  \item[Core or free foundation.] Creates the customer-owned AION brain and provides useful local
        intelligence without a mandatory paid model provider.
  \item[Boardroom.] Adds the principal business-management, departmental and connected-system
        capabilities.
  \item[Growth modules.] Add project portfolios, teams, budgets, divisions, locations, multiple
        entities and consolidated operation when required.
  \item[Sovereign or enterprise modules.] Add enterprise identity, single sign-on, private cloud,
        residency controls, advanced security, audit, operational assurance and contracted support.
\end{description}

Pricing and activation should consequently follow selected capabilities, verified usage and support
requirements rather than employee count.  A regulated sole practitioner may select Sovereign
controls, while a larger but operationally straightforward organisation may require only the
Boardroom and selected Growth modules.

\subsection{Onboarding Output Contract}

The requirements conversation produces a customer-reviewable record containing:

\begin{itemize}
  \item the facts supplied by the customer and their provenance;
  \item unresolved questions and assumptions;
  \item the derived operating-complexity components;
  \item recommended modules and the reason for each recommendation;
  \item identity, authority, residency, retention and approval implications;
  \item expected local or cloud compute requirements;
  \item optional commercial entitlements, kept separate from technical need;
  \item the minimum interface initially presented to each authorized role; and
  \item an exact approval record for activating the configuration.
\end{itemize}

The output is written into the customer's governed Business Map only after review.  Inferred facts
remain distinguishable from confirmed facts.  No model, installer or commercial package may silently
grant a user authority, connect an external service, move data to a provider or activate a paid
capability.

\subsection{Relationship to the Medium-Business Operating Foundation}

The implemented operating foundation already supports multi-entity groups, arbitrary divisions,
locations, matrix teams, project portfolios, currency-separated budgets, enterprise identity
subjects, delegated approvals, separation of duty and regional routing.  These capabilities are
reusable for organizations of any nominal size.  The requirements qualification determines which
of them become visible and configured for a particular customer.

The resulting rule is:

\begin{quote}
\textbf{AION qualifies operational requirements, progressively reveals the necessary capabilities,
and prices chosen service and support---it does not classify the customer by an unreliable size label.}
\end{quote}
